\documentclass[UTF8,zihao=false,12pt]{ctexart}
\usepackage[
  paperwidth=108mm,
  paperheight=192mm,
  top=10mm,
  bottom=10mm,
  left=9mm,
  right=9mm
]{geometry}
\usepackage{xcolor}
\usepackage{enumitem}
\usepackage{titlesec}
\usepackage{setspace}
\usepackage{fontspec}

\setmainfont{Times New Roman}
\setCJKmainfont{Songti SC}
\setlength{\parindent}{0pt}
\setlength{\parskip}{0.58em}
\setstretch{1.18}
\emergencystretch=2em
\pagestyle{empty}

\definecolor{titleblue}{HTML}{1F4E79}
\definecolor{muted}{HTML}{666666}

\titleformat{\section}
  {\Large\bfseries\color{titleblue}}
  {}{0pt}{}
\titleformat{\subsection}
  {\large\bfseries}
  {}{0pt}{}
\titlespacing*{\section}{0pt}{0.4em}{0.45em}
\titlespacing*{\subsection}{0pt}{0.45em}{0.25em}
\setlist[itemize]{leftmargin=1.2em,itemsep=0.25em,topsep=0.15em}

\newcommand{\slide}[2]{%
  \section*{#1}
  {\color{muted}#2}\par
  \vspace{0.4em}
}

\begin{document}

\begin{center}
  {\Huge\bfseries 数值实验与政策仿真}\\[0.7em]
  {\Large 3分钟手机讲稿}\\[0.8em]
  {\color{muted}逻辑主线：为什么做实验、实验说明了什么、政策含义是什么}
\end{center}

\vspace{1em}

这一部分主要回答一个问题：

\textbf{在港口设备电动化转型中，企业应该如何在需求不确定和政策补贴下做投资决策。}

我们做三组实验，分别回答三个问题：

\begin{itemize}
  \item 随机规划有没有必要；
  \item 补贴政策是否真的能推动减排；
  \item 当需求波动变强时，前面的结论是否仍然成立。
\end{itemize}

\newpage

\slide{第21页：总起}{五、数值实验与政策仿真}

这一页只做过渡。

可以这样说：

\textbf{接下来我汇报数值实验与政策仿真部分。}

前面的模型已经给出了随机规划框架，但仅有模型还不够。我们还需要通过实验说明：为什么要考虑需求不确定性，以及政策补贴在什么情况下能够真正产生作用。

\newpage

\slide{第22页：实验设计}{三组实验的 motivation}

我们设计了三组实验。

\textbf{第一组比较 DEP 和 EV。}

DEP 是考虑多个需求情景的随机规划模型；EV 是只看平均需求的模型。

这组实验的动机是：如果企业只按照平均需求规划，会不会低估高峰需求下的 EYC 配置需求？

\textbf{第二组改变补贴比例 $k$，从 0 到 1。}

动机是：补贴降低了企业投资成本，但它是否一定会带来更多 EYC 投资和更低排放？

\textbf{第三组改变需求波动水平 CV。}

动机是：现实需求波动不固定，所以要看在不同波动强度下，企业投资行为和补贴效果是否稳定。

\newpage

\slide{第23页：需求不确定性影响}{DEP vs EV}

第一组实验说明，随机规划确实有价值。

以中规模实例为例，EV 下 EYC 投入是 \textbf{11.20}，DEP 下提高到 \textbf{16.19}。

这说明如果只看平均需求，企业会低估电动化设备需求。

原因是 EV 把需求波动平滑掉了，看不到高峰情景；DEP 会考虑多个情景，所以会提前配置更多 EYC 作为缓冲。

大规模实例中，这个差距更明显，DEP 比 EV 多配置 \textbf{7.32} 的 EYC。

同时 VSS 为正，说明随机规划不是单纯增加模型复杂度，而是能降低不确定环境下的期望成本。

\textbf{这一页的结论：需求不确定性会推动企业预留更多 EYC。EYC 不只是替代柴油设备，也是一种应对高峰风险的保险性产能。}

\newpage

\slide{第24页：补贴灵敏度}{改变补贴比例 $k$}

第二组实验看补贴比例 $k$。

直觉上，补贴越高，企业越愿意投资 EYC，排放应该下降。

但结果显示，随着 $k$ 从 0 增加到 1，企业总成本确实下降，政府补贴支出增加；但总排放在大部分区间没有明显下降。

这说明当前参数下，补贴更多是在降低企业成本，而不是直接形成新增减排。

原因是基准情况下，企业已经完成了主要电动化投资和区域改造，所以继续提高补贴，不会明显改变实际运营中 EYC 和 DYC 的使用结构。

我们还设置了政策临界点：总排放相比 $k=0$ 下降至少 0.5\%。但结果显示没有明显临界点。

\textbf{这一页的结论：单一投资补贴可以降低企业负担，但不是充分的减排政策。真正减排还需要影响运营端，比如碳价、排放配额，或者和实际减排挂钩的补贴。}

\newpage

\slide{第25页：CV 扩展}{改变需求波动水平}

第三组实验改变需求波动水平 CV。

结果显示，CV 越高，企业越倾向于配置更多 EYC。

比如没有补贴时，CV=0.05 下 EYC 是 \textbf{12.63}，CV=0.30 时提高到 \textbf{22.74}。

这说明需求越不稳定，企业越需要额外 EYC 来应对高峰情景。

但同时，即使补贴提高带来了更多 EYC，排放也不一定同步下降。

这说明排放不只取决于买了多少 EYC，还取决于实际运营中是否用 EYC 替代了 DYC。

\textbf{这一页的结论：高不确定性会放大企业对 EYC 缓冲能力的需求，但补贴要真正转化为减排，还需要配合运营侧约束。}

\newpage

\slide{第26页：管理洞察}{总结和政策建议}

综合三组实验，我们有三个结论。

\textbf{第一，从企业角度，}需求波动越强，越不能只按平均需求规划。随机规划能帮助企业提前识别高峰风险，配置更稳健的 EYC。

\textbf{第二，从政府角度，}投资补贴主要是降低企业电动化成本。但如果企业在低补贴下已经完成主要投资，继续提高补贴的边际减排效果会变弱。

\textbf{第三，从政策设计角度，}补贴最好不要单独使用。更有效的是把投资补贴和碳交易价格、排放配额递减、实际减排绩效考核结合起来。

\vspace{0.4em}
\textbf{最后一句：}

我们的实验说明，随机规划能提升企业在不确定需求下的投资稳健性；而政策层面，单一补贴更多是成本分担工具，真正减排需要投资端激励和运营端碳约束协同发挥作用。

\newpage

\section*{可能被问到的问题（一）}
\subsection*{1. DEP 和 EV 的本质区别是什么？}

EV 只用平均需求做一个确定性规划；DEP 同时考虑多个需求情景和对应概率。

所以 EV 更像“按平时水平做计划”，DEP 更像“把高峰和低谷都提前纳入计划”。

\subsection*{2. 为什么要比较 DEP 和 EV？}

因为我们想证明随机规划有没有必要。

如果 DEP 和 EV 结果差不多，说明平均需求规划已经足够；但实验中 DEP 明显配置更多 EYC，且 VSS 为正，说明忽略不确定性会低估高峰风险。

\subsection*{3. 为什么 DEP 会配置更多 EYC？}

因为 EYC 购买和区域改造是一阶段决策，要在真实需求发生前确定。

DEP 会看到高需求情景下的服务压力，所以会提前保留更多 EYC 作为缓冲能力；EV 把需求平均掉，因此看不到这部分风险。

\subsection*{4. VSS 是什么？怎么解释？}

VSS 是随机规划价值，可以理解为“使用随机规划相比使用均值方案节省的期望成本”。

我们是最小化成本，所以 VSS 为正表示随机规划方案更好，也说明考虑不确定性确实有经济价值。

\subsection*{5. 如果老师问一句话总结第一组实验？}

可以说：

\textbf{第一组实验说明，只按平均需求规划会低估高峰风险；随机规划会通过增加 EYC 缓冲能力，获得更稳健的投资方案。}

\newpage

\section*{可能被问到的问题（二）}

\subsection*{6. 为什么补贴提高了，排放没有明显下降？}

因为在当前参数下，企业在低补贴时已经完成了主要电动化投资和区域改造。继续提高补贴主要降低企业成本，但不一定改变实际运营中 EYC 和 DYC 的使用比例。

\subsection*{7. 那是不是说明补贴没有用？}

不是。补贴有用，但主要作用是降低企业投资负担。若目标是减排，还需要碳价、排放配额或绩效型补贴，把政策作用传导到运营端。

\subsection*{8. 为什么 EYC 增加但排放不一定下降？}

因为排放取决于实际作业中 EYC 是否替代了 DYC，而不只取决于企业拥有多少 EYC。

也就是说，“投资端电动化能力增加”和“运营端实际减排”之间还需要调度和政策约束来连接。

\subsection*{9. 政策临界点为什么设成 0.5\%？}

0.5\% 是一个用于判定“排放是否出现明显下降”的经验阈值。

这样可以避免把非常小的数值波动误判成政策效果。我们的结果没有达到这个阈值，所以结论是：当前参数下没有形成显著减排临界点。

\subsection*{10. 如果老师问一句话总结第二组实验？}

可以说：

\textbf{第二组实验说明，单一投资补贴能降低企业成本，但不一定自动带来减排；要减排，需要政策影响实际运营排放结构。}

\newpage

\section*{可能被问到的问题（三）}

\subsection*{11. CV 表示什么？为什么改变 CV？}

CV 是需求波动水平。CV 越大，说明需求越不稳定，高峰需求风险越明显。

改变 CV 是为了检验结论是否稳健：如果需求波动更强，企业是否会更倾向于预留 EYC，补贴效果是否会改变。

\subsection*{12. 为什么 CV 越高，EYC 投入越多？}

因为高 CV 意味着高峰情景更突出。企业为了避免在高峰情景下依赖 DYC 或产生高运营成本，会提前配置更多 EYC。

所以在高不确定性下，EYC 更像一种“保险性产能”。

\subsection*{13. 补贴和碳交易价格为什么要协同？}

补贴主要作用在投资端，降低购买或改造 EYC 的成本。

碳价主要作用在运营端，提高使用 DYC 的排放成本。

两者协同后，企业既有动力投资 EYC，也有动力在实际运营中使用 EYC 替代 DYC。

\subsection*{14. 这部分实验有什么局限？}

第一，算例参数是基于模型设定生成的，不是完整真实港口数据。

第二，补贴效果和碳价、设备成本、排放因子有关，换一组参数后具体数值可能变化。

但我们的重点不是某一个数值，而是机制：需求不确定性会提高 EYC 缓冲需求，单一补贴未必充分减排。

\subsection*{15. 如果老师问你的核心贡献是什么？}

可以说：

\textbf{我们用实验把随机规划和政策仿真连接起来，说明了企业投资决策受需求风险影响，而政府政策要真正减排，不能只补贴投资端，还要约束运营端排放。}

\end{document}
